\section*{Capítulo \ref{ch:g7:identifying-substances} --- Identificar sustancias}

\begin{solution}{exo:g7:identifying-substances:1}
El ensayo del agua de cal: se hace burbujear el gas a través de agua de
cal transparente, que se enturbia si el gas es dióxido de carbono.
\end{solution}

\begin{solution}{exo:g7:identifying-substances:2}
El polvo es sulfato de cobre anhidro, y el líquido contiene agua.
\end{solution}

\begin{solution}{exo:g7:identifying-substances:3}
Oxígeno.
\end{solution}

\begin{solution}{exo:g7:identifying-substances:4}
El lápiz no se \omterm{def:g2:mixing-and-dissolving:dissolve}{disuelve} en el \omterm{def:g6:solutions-and-solubility:solute}{disolvente}. Una línea de tinta sería
arrastrada hacia arriba y separada, y estropearía el \omterm{def:g7:identifying-substances:chromatography}{cromatograma}.
\end{solution}

\begin{solution}{exo:g7:identifying-substances:5}
Si la mancha quedara por debajo del \omterm{def:g6:solutions-and-solubility:solute}{disolvente}, los colorantes se
disolverían en el \omterm{def:g6:solutions-and-solubility:solute}{disolvente} del vaso en lugar de subir por el papel.
\end{solution}

\begin{solution}{exo:g7:identifying-substances:6}
No, es una \omterm{def:g6:pure-substances-and-mixtures:mixture}{mezcla} de dos colorantes: probablemente el colorante
amarillo y el colorante azul de las referencias.
\end{solution}

\begin{solution}{exo:g7:identifying-substances:7}
Recoger un poco del gas en un tubo de ensayo y acercar una astilla
encendida a su boca: un pop agudo indica hidrógeno.
\end{solution}

\begin{solution}{exo:g7:identifying-substances:8}
El \omterm{def:g4:air-a-mixture-of-gases:air}{aire} espirado contiene mucho más dióxido de carbono que el \omterm{def:g4:air-a-mixture-of-gases:air}{aire}
fresco.
\end{solution}

\begin{solution}{exo:g7:identifying-substances:9}
El ensayo muestra que hay agua, no que esté sola: el agua salada, un
zumo o cualquier \omterm{def:g6:pure-substances-and-mixtures:mixture}{mezcla} acuosa también volverían azul el polvo.
\end{solution}

\begin{solution}{exo:g7:identifying-substances:10}
Tres colorantes (tres manchas). El colorante amarillo es el que subió
más alto.
\end{solution}

\begin{solution}{exo:g7:identifying-substances:11}
Por ejemplo: una astilla incandescente en cada frasco; el que la reaviva
contiene oxígeno. En los otros dos, una astilla encendida en la boca: un
pop agudo indica hidrógeno. El último frasco contiene dióxido de carbono
(se puede confirmar con agua de cal). Bastan dos o tres ensayos.
\end{solution}

\begin{solution}{exo:g7:identifying-substances:12}
Puede que el gas no contenga dióxido de carbono, o que contenga
demasiado poco para notarse en ese tiempo. Un ensayo de control con \omterm{def:g4:air-a-mixture-of-gases:air}{aire}
espirado, que se sabe que contiene dióxido de carbono, comprueba que el
agua de cal funciona.
\end{solution}

\begin{solution}{pb:g7:identifying-substances:1}
\textbf{1.} Para que las condiciones (papel, \omterm{def:g6:solutions-and-solubility:solute}{disolvente}, tiempo) sean
las mismas para las cuatro: solo así se pueden comparar las alturas de
las manchas.

\textbf{2.} Una \omterm{def:g6:pure-substances-and-mixtures:mixture}{mezcla}: su tinta da tres manchas separadas, tres
colorantes.

\textbf{3.} Dos colorantes. No: la tinta de la nota tiene tres
colorantes, y el rotulador 1 no tiene ni el colorante rojo ni el
amarillo.

\textbf{4.} Sus tres manchas no están a las mismas alturas que las de la
nota: son tres colorantes distintos.

\textbf{5.} El rotulador 2: sus tres manchas están exactamente a las
alturas de las tres manchas de la nota.

\textbf{6.} Oxígeno.

\textbf{7.} Hidrógeno.

\textbf{8.} Dióxido de carbono.

\textbf{9.} Un ensayo en blanco (un control). Muestra que el agua de cal
no se enturbia sola, ni con \omterm{def:g4:air-a-mixture-of-gases:air}{aire} normal en ese tiempo: así que la
turbidez con el gas Z se debe de verdad al dióxido de carbono.

\textbf{10.} ``Agua --- sulfato de cobre anhidro --- de blanco pasa a
azul'': el líquido contiene agua.

\textbf{11.} No. Un ensayo positivo muestra que hay agua, no que no haya
nada más: el líquido podría ser una disolución u otra \omterm{def:g6:pure-substances-and-mixtures:mixture}{mezcla} acuosa.

\textbf{12.} La tinta de la nota contiene 3 colorantes, y fue el
rotulador 2 el que escribió la nota.
\end{solution}
